\section{Clase 27 Gabriel}

\subsection*{El haz dualizante de $\mathbb{P}^1_K$}

El otro día vimos el teorema de representabilidad aplicado al funtor
\[
F:\operatorname{Coh}(X)^{\operatorname{op}}\longrightarrow K\text{-Vect},
\qquad
\mathcal{M}\longmapsto H^1(X,\mathcal{M})^\vee.
\]
Este funtor es representable. Es decir, existe un haz cuasicoherente
\[
\omega
\]
tal que, para todo haz coherente $\mathcal{M}$, se tiene un isomorfismo natural
\[
F(\mathcal{M}) = H^1(X,\mathcal{M})^\vee
\simeq
\operatorname{Hom}_{\mathcal{O}_X}(\mathcal{M},\omega).
\]
El haz $\omega$ se llama el \textbf{haz dualizante} de $X$.

La pregunta que queremos responder ahora es:

\[
\text{¿quién es } \omega \text{ cuando } X=\mathbb{P}^1_K\text{?}
\]

\begin{theorem}
Si
\[
X=\mathbb{P}^1_K,
\]
entonces
\[
\omega\simeq \mathcal{O}(-2).
\]
\end{theorem}

Recordemos que
\[
\mathcal{O}(-2)
\]
puede pensarse como el haz de funciones racionales que tienen al menos un cero de orden $2$ en el punto del infinito.

Además, ya calculamos la cohomología de los haces $\mathcal{O}(n)$ en $\mathbb{P}^1_K$:
\[
\dim_K H^p(\mathbb{P}^1_K,\mathcal{O}(n))=
\begin{cases}
n+1, & p=0,\ n\ge 0,\\
0, & p=0,\ n<0,\\
-n-1, & p=1,\ n\le -2,\\
0, & p=1,\ n\ge -1,\\
0, & p\neq 0,1.
\end{cases}
\]
En particular,
\[
\dim_K H^1(\mathbb{P}^1_K,\mathcal{O}(-2))=1.
\]
Tomamos entonces un elemento no nulo
\[
\xi\in H^1(\mathbb{P}^1_K,\mathcal{O}(-2))^\vee.
\]

\begin{lemma}
La pareja
\[
(\mathcal{O}(-2),\xi)
\]
es mínima.
\end{lemma}

\begin{proof}
Supongamos que existe un morfismo epiyectivo de parejas
\[
(\mathcal{O}(-2),\xi)\longrightarrow (\mathcal{M},\xi_{\mathcal{M}}).
\]
Esto significa que tenemos un epimorfismo de haces
\[
\varphi:\mathcal{O}(-2)\longrightarrow \mathcal{M}
\]
tal que
\[
F(\varphi)(\xi_{\mathcal{M}})=\xi.
\]

Si $\ker\varphi$ tuviese rango $1$, entonces el cociente $\mathcal{M}$ sería un haz de torsión. Pero todo haz coherente de torsión en $\mathbb{P}^1_K$ tiene soporte en un número finito de puntos cerrados, luego es una suma finita de haces rascacielos. En particular es acíclico, y por tanto
\[
H^1(\mathbb{P}^1_K,\mathcal{M})=0.
\]
Esto implica
\[
\xi_{\mathcal{M}}=0,
\]
y entonces
\[
\xi=F(\varphi)(\xi_{\mathcal{M}})=0,
\]
contradicción.

Por tanto, $\ker\varphi$ no puede tener rango $1$. Como $\mathcal{O}(-2)$ es localmente libre de rango $1$, su núcleo sólo puede tener rango $0$. Además, al ser un subhaz de un haz localmente libre, no puede ser de torsión salvo que sea cero. Luego
\[
\ker\varphi=0.
\]
Así, $\varphi$ es a la vez mono y epi, luego es un isomorfismo. Esto demuestra que la pareja $(\mathcal{O}(-2),\xi)$ es mínima.
\end{proof}

Por el teorema de representabilidad, el haz $\omega$ representa el funtor $F$. La clase $\xi$ determina entonces un morfismo
\[
\mathcal{O}(-2)\longrightarrow \omega.
\]
Por la minimalidad de la pareja $(\mathcal{O}(-2),\xi)$, este morfismo identifica $\mathcal{O}(-2)$ con un subhaz de $\omega$. Por tanto tenemos una sucesión exacta
\[
0\longrightarrow \mathcal{O}(-2)
\longrightarrow \omega
\longrightarrow \mathcal{K}
\longrightarrow 0.
\]
Para concluir que
\[
\omega\simeq \mathcal{O}(-2),
\]
basta probar que
\[
\mathcal{K}=0.
\]

\subsection*{Una notación}

Si $\mathcal{M}$ es un haz cuasicoherente en $\mathbb{P}^1_K$, escribiremos
\[
\mathcal{M}(n)
=
\mathcal{M}\otimes_{\mathcal{O}_{\mathbb{P}^1_K}}\mathcal{O}(n).
\]
Tensorizando la sucesión exacta anterior por $\mathcal{O}(n)$, obtenemos
\[
0\longrightarrow \mathcal{O}(n-2)
\longrightarrow \omega(n)
\longrightarrow \mathcal{K}(n)
\longrightarrow 0.
\]
La sucesión sigue siendo exacta porque $\mathcal{O}(n)$ es localmente libre.

\begin{lemma}
Para todo $n\in\mathbb{Z}$, se tiene un isomorfismo natural
\[
\mathcal{O}(n)^\vee\simeq \mathcal{O}(-n).
\]
\end{lemma}

\begin{proof}
Recordemos que
\[
\mathcal{O}(n)^\vee=
\mathcal{H}om_{\mathcal{O}_{\mathbb{P}^1_K}}(\mathcal{O}(n),\mathcal{O}_{\mathbb{P}^1_K}).
\]
Es una cuestión local.

Sea $U\subset \mathbb{P}^1_K$ un abierto. Si $\infty\notin U$, entonces
\[
\mathcal{O}(n)|_U\simeq \mathcal{O}_U,
\]
y por tanto
\[
\mathcal{O}(n)^\vee(U)
\simeq
\mathcal{O}(U)
\simeq
\mathcal{O}(-n)(U).
\]

Supongamos ahora que $\infty\in U$. Tomamos como parámetro local en $\infty$
\[
\pi=\frac{1}{t}.
\]
Entonces, localmente,
\[
\mathcal{O}(n)(U)=\pi^{-n}\mathcal{O}(U).
\]
Por tanto,
\[
\mathcal{O}(n)^\vee(U)=
\operatorname{Hom}_{\mathcal{O}(U)}
(\pi^{-n}\mathcal{O}(U),\mathcal{O}(U)).
\]
Un homomorfismo queda determinado por la imagen de $\pi^{-n}$, y esta imagen debe ser un elemento $g\in\mathcal{O}(U)$. Equivalentemente, podemos identificar estos homomorfismos con los elementos $h\in K(t)$ tales que
\[
h\pi^{-n}\in \mathcal{O}(U).
\]
Es decir,
\[
h\in \pi^n\mathcal{O}(U).
\]
Pero
\[
\pi^n\mathcal{O}(U)=\mathcal{O}(-n)(U).
\]
Así,
\[
\mathcal{O}(n)^\vee(U)\simeq \mathcal{O}(-n)(U).
\]
Estos isomorfismos son compatibles con las restricciones, luego definen un isomorfismo de haces
\[
\mathcal{O}(n)^\vee\simeq \mathcal{O}(-n).
\]
\end{proof}

\begin{lemma}
Para todo $n\ge 1$, se tiene
\[
H^0(\mathbb{P}^1_K,\mathcal{K}(n))=0.
\]
\end{lemma}

\begin{proof}
Primero observamos que, para todo haz cuasicoherente $\mathcal{M}$, hay una identificación natural
\[
\mathcal{M}(X)=
\operatorname{Hom}_{\mathcal{O}_X}(\mathcal{O}_X,\mathcal{M}).
\]
En efecto, a una sección $m\in \mathcal{M}(X)$ le asociamos el morfismo
\[
\mathcal{O}_X\longrightarrow \mathcal{M},
\qquad
f\longmapsto fm.
\]
Recíprocamente, todo morfismo $\mathcal{O}_X\to \mathcal{M}$ queda determinado por la imagen de $1$.

Aplicando esto a $\omega(n)$, obtenemos
\[
H^0(\mathbb{P}^1_K,\omega(n))=
\operatorname{Hom}_{\mathcal{O}_{\mathbb{P}^1_K}}(\mathcal{O},\omega(n)).
\]
Como
\[
\omega(n)=\omega\otimes\mathcal{O}(n),
\]
y como
\[
\mathcal{O}(n)^\vee\simeq \mathcal{O}(-n),
\]
tenemos
\[
\operatorname{Hom}_{\mathcal{O}_{\mathbb{P}^1_K}}(\mathcal{O},\omega(n))
\simeq
\operatorname{Hom}_{\mathcal{O}_{\mathbb{P}^1_K}}(\mathcal{O}(-n),\omega).
\]
Por representabilidad de $\omega$,
\[
\operatorname{Hom}_{\mathcal{O}_{\mathbb{P}^1_K}}(\mathcal{O}(-n),\omega)
\simeq
H^1(\mathbb{P}^1_K,\mathcal{O}(-n))^\vee.
\]
Luego
\[
\dim_K H^0(\mathbb{P}^1_K,\omega(n))
=
\dim_K H^1(\mathbb{P}^1_K,\mathcal{O}(-n)).
\]
Para $n\ge 1$, esto vale
\[
\dim_K H^0(\mathbb{P}^1_K,\omega(n))=n-1.
\]

Por otro lado,
\[
\dim_K H^0(\mathbb{P}^1_K,\mathcal{O}(n-2))=n-1
\]
para todo $n\ge 1$.

Tomamos ahora cohomología en la sucesión exacta
\[
0\longrightarrow \mathcal{O}(n-2)
\longrightarrow \omega(n)
\longrightarrow \mathcal{K}(n)
\longrightarrow 0.
\]
Obtenemos
\[
0\longrightarrow
H^0(\mathbb{P}^1_K,\mathcal{O}(n-2))
\longrightarrow
H^0(\mathbb{P}^1_K,\omega(n))
\longrightarrow
H^0(\mathbb{P}^1_K,\mathcal{K}(n))
\longrightarrow
H^1(\mathbb{P}^1_K,\mathcal{O}(n-2)).
\]
Pero para $n\ge 1$,
\[
H^1(\mathbb{P}^1_K,\mathcal{O}(n-2))=0.
\]
Así, tenemos una sucesión exacta
\[
0\longrightarrow
H^0(\mathbb{P}^1_K,\mathcal{O}(n-2))
\longrightarrow
H^0(\mathbb{P}^1_K,\omega(n))
\longrightarrow
H^0(\mathbb{P}^1_K,\mathcal{K}(n))
\longrightarrow 0.
\]
Los dos primeros espacios tienen la misma dimensión, igual a $n-1$. Por tanto el primer morfismo es un isomorfismo y
\[
H^0(\mathbb{P}^1_K,\mathcal{K}(n))=0.
\]
\end{proof}

Ya podemos terminar la demostración del teorema.

\begin{proof}[Conclusión de la demostración del teorema]
Queremos probar que
\[
\mathcal{K}=0.
\]

Primero veamos que $\mathcal{K}$ no tiene torsión. Si tuviese torsión no nula, esa parte de torsión tendría soporte en un número finito de puntos cerrados. Los haces rascacielos siempre tienen secciones globales no nulas. Además, tensorizar por $\mathcal{O}(n)$ no cambia el hecho de tener soporte finito. Por tanto, para todo $n$, el haz $\mathcal{K}(n)$ tendría alguna sección global no nula. Esto contradice que
\[
H^0(\mathbb{P}^1_K,\mathcal{K}(n))=0
\]
para $n\ge 1$. Luego $\mathcal{K}$ no tiene torsión.

Nos queda ver que tampoco tiene parte libre. Basta probar que su tallo genérico es cero:
\[
\mathcal{K}_\eta=0.
\]
Supongamos, por contradicción, que existe
\[
0\neq s\in \mathcal{K}_\eta.
\]
Entonces $s$ está definido en algún abierto no vacío
\[
U\subset \mathbb{P}^1_K.
\]
El complemento
\[
\mathbb{P}^1_K\setminus U
\]
es finito. Escribamos
\[
\mathbb{P}^1_K\setminus U=\{p_1,\dots,p_r\}.
\]
La sección racional $s$ puede tener polos en esos puntos. Tomamos enteros $n_i$ suficientemente grandes para compensar esos polos, y consideramos el divisor efectivo
\[
D=\sum_{i=1}^r n_i p_i.
\]
Entonces $s$ define una sección global no nula de
\[
\mathcal{K}\otimes L_D.
\]
Como en $\mathbb{P}^1_K$ se tiene
\[
\operatorname{Pic}(\mathbb{P}^1_K)\simeq \mathbb{Z},
\]
y
\[
L_D\simeq \mathcal{O}(\deg D),
\]
obtenemos una sección global no nula de
\[
\mathcal{K}(\deg D).
\]
Tomando $\deg D\ge 1$, esto contradice el lema anterior, que dice que
\[
H^0(\mathbb{P}^1_K,\mathcal{K}(n))=0
\qquad
\text{para todo } n\ge 1.
\]
Por tanto
\[
\mathcal{K}_\eta=0.
\]

Hemos probado que $\mathcal{K}$ no tiene torsión y que su tallo genérico es cero. En consecuencia,
\[
\mathcal{K}=0.
\]
Así, la sucesión exacta
\[
0\longrightarrow \mathcal{O}(-2)
\longrightarrow \omega
\longrightarrow \mathcal{K}
\longrightarrow 0
\]
se reduce a un isomorfismo
\[
\omega\simeq \mathcal{O}(-2).
\]
\end{proof}

\begin{corollary}[Dualidad de Serre en $\mathbb{P}^1_K$]
Para todo haz coherente $\mathcal{M}$ en $\mathbb{P}^1_K$, se tiene un isomorfismo natural
\[
H^1(\mathbb{P}^1_K,\mathcal{M})^\vee
\simeq
\operatorname{Hom}_{\mathcal{O}_{\mathbb{P}^1_K}}(\mathcal{M},\mathcal{O}(-2)).
\]
\end{corollary}

\begin{remark}
La idea geométrica es que el haz dualizante de $\mathbb{P}^1_K$ es el haz que representa los funcionales lineales sobre la cohomología $H^1$. El cálculo anterior muestra que dicho representante es precisamente $\mathcal{O}(-2)$, que coincide con el haz canónico de la recta proyectiva.
\end{remark}
